\ifdefined\gkver\else\def\gkver{student}\fi
\documentclass[\gkver]{gaokaozhenti}
\begin{document}

\chapter{1983年全国}

\begin{enumerate}
\item
%% number: 1
%% paperName: 1983年全国高考第 1 题 4 分
%% typeId: 3
%% type: 填空题
%% chapter: 运动和力
%% point: 力学单位制
%% method: 无
%% score: 4
%% degree: 850
%% duplicateId: 0
%% body:
写出下列物理量的单位或数值：

北京地区的重力加速度是 $9.801$\tkanswer{$\Umsq$}；冰的熔点是\tkanswer{273}$\UK$；第一宇宙速度大约是\tkanswer{7.9} 千米/秒；基本电荷 $e = 1.60\times10^{-19}$\tkanswer{$\UC$}。

%% answer: 米/秒^2；273；7.9；库仑
\jdanswer{
$\Umsq$，273，7.9，$\UC$
}

%% memo:
\memoanswer{
重力加速度的单位是 $\Umsq$；冰的熔点是 $0\UdC$，由 $T=273+t$ 得 $273\UK$；第一宇宙速度大约是 $7.9\Ukms$；基本电荷 $e=1.60\times10^{-19}\UC$。

评分标准：全题 20 分，每小题 4 分。各题中，在每一横线上空白处填入正确答案的，按下列该答案后面方括号内的分数给分；错误的，给 0 分。米/秒$^{2}$（1分），273（1分），7.9（1分。答 7.8 或 8 的，同样给分），库仑（1分）。
}

\item
%% number: 2
%% paperName: 1983年全国高考第 2 题 4 分
%% typeId: 3
%% type: 填空题
%% chapter: 曲线运动
%% point: 斜抛运动
%% method: 无
%% score: 4
%% degree: 650
%% duplicateId: 0
%% body:
以 $60^{\circ}$ 的仰角抛出一物体（空气阻力不计）。它到达最高点时的动量的大小跟刚抛出时的动量的大小之比等于\tkanswer{$\frac{1}{2}$}。物体到达最高点时的动能跟刚抛出时的动能之比等于\tkanswer{$\frac{1}{4}$}。

%% answer: 1/2；1/4
\jdanswer{
$\frac{1}{2}$，$\frac{1}{4}$
}

%% memo:
\memoanswer{
由题意知 $\frac{p_{2}}{p_{1}}=\frac{v_{2}}{v_{1}}=\frac{v_{0}\cos60^{\circ}}{v_{0}}=\frac{1}{2}$，$\frac{E_{2}}{E_{1}}=\frac{v_{2}^{2}}{v_{1}^{2}}=\frac{v_{0}^{2}\cos^{2}60^{\circ}}{v_{0}^{2}}=\frac{1}{4}$。

评分标准：全题 20 分，每小题 4 分。
}

\item
%% number: 3
%% paperName: 1983年全国高考第 3 题 4 分
%% typeId: 3
%% type: 填空题
%% chapter: 光学
%% point: 光电效应
%% method: 无
%% score: 4
%% degree: 650
%% duplicateId: 0
%% body:
使金属钠产生光电效应的光的最长波长是 5000 埃。因此，金属钠的逸出功 $W =$\tkanswer{$3.98\times10^{-19}$}$\UJ$。现在用频率在 $3.90\times10^{14}\UHz$ 到 $7.50\times10^{14}\UHz$ 范围内的光照射钠，那么，使钠产生光电效应的频率范围是从\tkanswer{$6.00\times10^{14}$}$\UHz$ 到\tkanswer{$7.50\times10^{14}$}$\UHz$。（普朗克恒量 $h = 6.63\times10^{-34}\UJcdots$）

%% answer: 3.98e-19；6.00e14；7.50e14
\jdanswer{
$3.98\times10^{-19}$，$6.00\times10^{14}$，$7.50\times10^{14}$
}

%% memo:
\memoanswer{
由光电效应知逸出功为
\[ W=h\frac{c}{\lambda_{0}}=6.63\times10^{-34}\times\frac{3\times10^{8}}{5000\times10^{-10}}\UJ=3.98\times10^{-19}\UJ; \]
使金属钠产生光电效应的光的最低频率为
\[ \nu=\frac{c}{\lambda}=\frac{3\times10^{8}}{5000\times10^{-10}}\UHz=6.00\times10^{14}\UHz, \]
所以使钠产生光电效应的频率范围是从 $6.00\times10^{14}\UHz$ 到 $7.50\times10^{14}\UHz$。

评分标准：全题 20 分，每小题 4 分。$3.98\times10^{-19}$（2分），$6.00\times10^{14}$，$7.50\times10^{14}$（2分）。本小题答两位有效数字也可。不填数值答案的，本答案给 0 分。$6.00\times10^{14}$ 和 $7.50\times10^{14}$ 两个数中有一个错的，本答案给 0 分。
}

\item
%% number: 4
%% paperName: 1983年全国高考第 4 题 4 分
%% typeId: 3
%% type: 填空题
%% chapter: 原子和原子核
%% point: 玻尔模型
%% method: 无
%% score: 4
%% degree: 750
%% duplicateId: 0
%% body:
已知氢原子的基态能级是 $E_{1} = -13.6\UeV$。第二能级是 $E_{2} = -3.4\UeV$。如果氢原子吸收\tkanswer{10.2}$\UeV$ 的能量，它即可由基态跃迁到第二能级。如果氢原子再获得 $1.89\UeV$ 的能量，它还可由第二能级跃迁到第三能级。因此，氢原子的第三能级 $E_{3} =$\tkanswer{$-1.51$}$\UeV$。

%% answer: 10.2；-1.51
\jdanswer{
10.2，$-1.51$
}

%% memo:
\memoanswer{
由玻尔理论知
\[ \Delta E=E_{2}-E_{1}=(-3.4\UeV)-(-13.6\UeV)=10.2\UeV, \]
\[ E_{3}=E_{2}+\Delta E'=(-3.4\UeV)+1.89\UeV=-1.51\UeV. \]

评分标准：全题 20 分，每小题 4 分。10.2（2分），$-1.51$（2分）。
}

\item
%% number: 5
%% paperName: 1983年全国高考第 5 题 4 分
%% typeId: 7
%% type: 作图题
%% chapter: 光学
%% point: 透镜
%% method: 无
%% score: 4
%% degree: 650
%% duplicateId: 0
%% body:
物体 AB 被置于一薄凸透镜的焦点 F 和光心 O 之间，并垂直于透镜的主轴。透镜的大小、焦点的位置、物体 AB 的长度和位置都如附图所示。

\begin{enumerate}
	\item
	在下图上用作图法（以透镜中通过光心的一段虚线代表薄透镜）画出凸透镜成像光路和像。

	\item
	要想看到物体 AB 的完整的像，眼睛必须处在某一范围内。试作图确定下图上的这一范围（用斜线标明）。
\end{enumerate}
\drawpicanswer{\onepicture[align=right, w=9cm]{figs/05.png}}{\onepicture[align=right, w=8cm]{figs/05_an.png}}

%% answer: （1）成像光路图如图所示；（2）见答案图
\jdanswer{
\begin{enumerate}
	\item
	成像光路图如图所示。

	\item
	眼睛位置在图中斜线标出的范围内。
\end{enumerate}
}

%% memo:
\memoanswer{
（1）成像光路图如图所示（光线 AM、AN、BM、BN 不属于要求的成像光路图）。只正确画出 A（或 B）的成像光路并按比例得出全像的，也给 2 分。光路图中，未标明光线进行方向或未画出全像的，扣 1 分。

（2）眼睛位置在图中斜线标出的范围内。只画出 $A'MP$ 和 $B'NP$ 两条线并用斜线标出范围的，也给 2 分。$A'M$ 和 $B'N$ 画成实线的，不扣分。作图不太准确的，不扣分。凡用斜线标出的范围与参考答案不符的，不给分。

评分标准：全题 20 分，每小题 4 分。
}

\item
%% number: 6
%% paperName: 1983年全国高考第 6 题 3 分
%% typeId: 1
%% type: 单选题
%% chapter: 电路
%% point: 闭合电路欧姆定律
%% method: 无
%% score: 3
%% degree: 650
%% duplicateId: 0
%% body:
在右图所示的电路里，$r$ 是电源的内阻，$R_{1}$ 和 $R_{2}$ 是外电路中的电阻。如果用 $P_{r}$、$P_{1}$ 和 $P_{2}$ 分别表示电阻 $r$、$R_{1}$ 和 $R_{2}$ 上所消耗的功率，当 $R_{1} = R_{2} = r$ 时，$P_{r}:P_{1}:P_{2}$ 等于\xzanswer{D}

\onepicture[w=5cm]{figs/06.png}

\fourchoices[answer=D]
{$1:1:1$}
{$2:1:1$}
{$1:4:4$}
{$4:1:1$}

%% answer: D
%% memo:
\memoanswer{
由题意知，电流关系 $I_{r} = 2I_{R_{1}} = 2I_{R_{2}}$，由 $P = I^{2}R$ 得，$P_{r} : P_{1} : P_{2} = 4 : 1 : 1$，故 D 正确。

评分说明：全题 15 分，每小题 3 分。每小题答案正确的给 3 分；答案错误的给 $-$1 分；未答的给 0 分。每小题选择了两个或两个以上答案的，无论答案写在括号内或括号旁，都给 $-$1 分。五个小题分数的代数和如果是正数或 0，即为本题全题得分；如果是负数，本题全题得分记作 0。
}

\item
%% number: 7
%% paperName: 1983年全国高考第 7 题 3 分
%% typeId: 1
%% type: 单选题
%% chapter: 运动和力
%% point: 动量守恒定律
%% method: 无
%% score: 3
%% degree: 750
%% duplicateId: 0
%% body:
一个静止的、质量为 $M$ 的不稳定原子核，当它放射出质量为 $m$，速度 $v$ 的粒子后，原子核的剩余部分的速度 $v'$ 等于\xzanswer{B}

\fourchoices[answer=B]
{$-v$}
{$-\frac{m}{M-m}v$}
{$-\frac{m}{m-M}v$}
{$-\frac{m}{M}v$}

%% answer: B
%% memo:
\memoanswer{
不稳定原子核放射时，由动量守恒得 $0 = mv + (M - m)u$，解得 $u = -\frac{mv}{M - m}$，故 B 正确。

评分说明：全题 15 分，每小题 3 分。
}

\item
%% number: 8
%% paperName: 1983年全国高考第 8 题 3 分
%% typeId: 1
%% type: 单选题
%% chapter: 电路
%% point: 电容
%% method: 无
%% score: 3
%% degree: 650
%% duplicateId: 0
%% body:
在右图所示的电路中，电容器 A 的电容 $C_{A} = 30\UuF$，电容器 B 的电容 $C_{B} = 10\UuF$。在开关 $K_{1}$、$K_{2}$ 都是断开的情况下，分别给电容器 A、B 充电。充电后，M 点的电势比 N 点高 $5\UV$，O 点的电势比 P 点低 $5\UV$。然后把 $K_{1}$、$K_{2}$ 都接通，接通后 M 点的电势比 N 点高\xzanswer{C}

\onepicture[w=6cm]{figs/08.png}

\fourchoices[answer=C]
{$10\UV$}
{$5\UV$}
{$2.5\UV$}
{$0\UV$}

%% answer: C
%% memo:
\memoanswer{
在开关 $K_{1}$、$K_{2}$ 都断开的情况下，由题意知，$U_{A}=5\UV$，$U_{B}=5\UV$，电容器 A 的电量 $Q_{A}=C_{A}U_{A}=150\UuC$，电容器 B 的电量 $Q_{B}=C_{B}U_{B}=50\UuC$，把 $K_{1}$、$K_{2}$ 都接通后，并联电容器的电量 $Q=Q_{A}+(-Q_{B})=100\UuC$，并联电容器的电容 $C=C_{A}+C_{B}=40\UuF$，所以并联电容器的电压 $U=\frac{Q}{C}=2.5\UV$，故 C 正确。

评分说明：全题 15 分，每小题 3 分。
}

\item
%% number: 9
%% paperName: 1983年全国高考第 9 题 3 分
%% typeId: 1
%% type: 单选题
%% chapter: 原子和原子核
%% point: 天然放射性
%% method: 无
%% score: 3
%% degree: 550
%% duplicateId: 0
%% body:
某原子核 A 的衰变过程如下（符号 $\xrightarrow{\beta}$ 表示放射一个 $\beta$ 粒子，$\xrightarrow{\alpha}$ 表示放射一个 $\alpha$ 粒子）：A $\xrightarrow{\beta}$ B $\xrightarrow{\alpha}$ C。\xzanswer{D}

\fourchoices[answer=D]
{核 A 的中子数减核 C 的中子数等于 2}
{核 A 的质量数减核 C 的质量数等于 5}
{原子核为 A 的中性原子中的电子数比原子核为 B 的中性原子中的电子数多 1}
{核 C 的质子数比核 A 的质子数少 1}

%% answer: D
%% memo:
\memoanswer{
原子核的衰变方程为 $\ce{^{M}_{Z}A -> ^{M}_{Z+1}B + ^{0}_{-1}e}$，$\ce{^{M}_{Z+1}B -> ^{M-4}_{Z-1}C + ^{4}_{2}He}$。核 A 的中子数减核 C 的中子数等于 $(M - Z) - [(M - 4) - (Z - 1)] = 3$，故 A 错误；核 A 的质量数减核 C 的质量数等于 $M - (M - 4) = 4$，故 B 错误；原子核为 A 的中性原子中的电子数（等于 $Z$）比原子核为 B 的中性原子中的电子数（等于 $Z + 1$）少 1，故 C 错误；核 C 的质子数（等于 $Z - 1$）比核 A 的质子数（等于 $Z$）少 1，故 D 正确。

评分说明：全题 15 分，每小题 3 分。
}

\item
%% number: 10
%% paperName: 1983年全国高考第 10 题 3 分
%% typeId: 1
%% type: 单选题
%% chapter: 电磁感应
%% point: 楞次定律
%% method: 无
%% score: 3
%% degree: 450
%% duplicateId: 0
%% body:
如下图所示，在两根平行长直导线 M、N 中，通过同方向同强度的电流。导线框 ABCD 和两导线在同一平面内。线框沿着与两导线垂直的方向，自右向左在两导线间匀速移动。在移动过程中，线框中感生电流的方向\xzanswer{B}

\onepicture[w=5cm]{figs/10.png}

\fourchoices[answer=B]
{沿 ABCDA，不变}
{沿 ADCBA，不变}
{由 ABCDA 变成 ADCBA}
{由 ADCBA 变成 ABCDA}

%% answer: B
%% memo:
\memoanswer{
在题图中，两导线间的磁场方向：中线以左向里，中线以右向外，磁感应强度为越靠近中线越小，中线上为 0。线框沿着与两导线垂直的方向，自右向左在两导线间匀速移动时，在中线以右，磁通量向外，并逐渐减少，由楞次定律知，感应电流的磁场方向向外，感应电流方向为 ADCBA；在中线以左，磁通量向里，并逐渐增大，由楞次定律知，感应电流的磁场方向向外，感应电流方向仍为 ADCBA，故 B 正确。

评分说明：全题 15 分，每小题 3 分。
}

\item
%% number: 11
%% paperName: 1983年全国高考第 11 题 3 分
%% typeId: 2
%% type: 多选题
%% chapter: 机械波
%% point: 波的图象
%% method: 无
%% score: 3
%% degree: 550
%% duplicateId: 0
%% body:
右图中的曲线是一列简谐横波在某一时刻的图象。根据这个图可以确定\xzanswer{BC}

\onepicture[w=7cm]{figs/11.png}

\fourchoices[answer=BC]
{周期}
{波长}
{振幅}
{波速}

%% answer: BC
%% memo:
\memoanswer{
由图知波长 $20\Ucm$，振幅 $5\Ucm$，故 BC 正确。

评分说明：全题 15 分，每小题 3 分。每小题答案全部选对的给 3 分；未选全或有选错的给 0 分。填写在方括号外的号码不作为考生的答案。
}

\item
%% number: 12
%% paperName: 1983年全国高考第 12 题 3 分
%% typeId: 2
%% type: 多选题
%% chapter: 光学
%% point: 透镜
%% method: 无
%% score: 3
%% degree: 850
%% duplicateId: 0
%% body:
以下光学器件是用凸透镜制成的，或相当于一个凸透镜。\xzanswer{ACD}

\fourchoices[answer=ACD]
{放大镜}
{近视眼镜}
{远视眼镜}
{幻灯机镜头}

%% answer: ACD
%% memo:
\memoanswer{
放大镜、远视眼镜和幻灯机镜头是凸透镜，故 ACD 正确。

评分说明：全题 15 分，每小题 3 分。
}

\item
%% number: 13
%% paperName: 1983年全国高考第 13 题 3 分
%% typeId: 2
%% type: 多选题
%% chapter: 电磁感应
%% point: 自感与互感，涡流
%% method: 无
%% score: 3
%% degree: 850
%% duplicateId: 0
%% body:
日光灯电路主要由镇流器、起动器和灯管组成。在日光灯正常工作的情况下\xzanswer{ABD}

\fourchoices[answer=ABD]
{灯管点燃发光后，起动器中两个触片是分离的}
{灯管点燃发光后，镇流器起降压限流作用}
{镇流器起整流作用}
{镇流器给日光灯的开始点燃提供瞬时高电压}

%% answer: ABD
%% memo:
\memoanswer{
日光灯开始启动时，起动器中两个触片受热而接触，然后电阻减小发热减少而分离，由于自感而产生瞬时高电压，启动日光灯管；灯管启动发光后，由于电感对交流电的阻碍作用，镇流器起降压限流作用，故 ABD 正确。

评分说明：全题 15 分，每小题 3 分。
}

\item
%% number: 14
%% paperName: 1983年全国高考第 14 题 3 分
%% typeId: 2
%% type: 多选题
%% chapter: 气体性质
%% point: 物体的内能
%% method: 无
%% score: 3
%% degree: 650
%% duplicateId: 0
%% body:
一个铜块和一个铁块，质量相等，铜块的温度 $T_{1}$ 比铁块的温度 $T_{2}$ 高。当它们接触在一起时，如果不和外界交换能量，那么\xzanswer{AD}

\fourchoices[answer=AD]
{从两者开始接触到热平衡的整个过程中，铜块放出的总热量等于铁块吸收的总热量}
{在两者达到热平衡以前的任一段时间内，铜块放出的热量不等于铁块吸收的热量}
{达到热平衡时，铜块的温度 $T = \frac{T_{1} + T_{2}}{2}$}
{达到热平衡时，两者的温度相等}

%% answer: AD
%% memo:
\memoanswer{
由热传递和热力学第零定律知，故 AD 正确、B 错误；铜与铁的比热容不同，所以平衡时的温度不等于 $\frac{T_{1}+T_{2}}{2}$，故 C 错误。

评分说明：全题 15 分，每小题 3 分。
}

\item
%% number: 15
%% paperName: 1983年全国高考第 15 题 3 分
%% typeId: 2
%% type: 多选题
%% chapter: 气体性质
%% point: 气体多过程
%% method: 图象法
%% score: 3
%% degree: 450
%% duplicateId: 0
%% body:
一定量的理想气体，处在某一初始状态。现在要使它的温度经过状态变化后回到初始状态的温度，用下列哪些过程可能实现？\xzanswer{AD}

\fourchoices[answer=AD]
{先保持压强不变而使它的体积膨胀，接着保持体积不变而减小压强}
{先保持压强不变而使它的体积减小，接着保持体积不变而减小压强}
{先保持体积不变而增大压强，接着保持压强不变而使它的体积膨胀}
{先保持体积不变而减小压强，接着保持压强不变而使它的体积膨胀}

%% answer: AD
%% memo:
\memoanswer{
气体等压膨胀时，吸热，即温度升高，接着等容降压，放热，温度降低，可回到初始状态，故 A 正确；气体等压减容，放热，温度降低，接着等容降压，放热，温度继续降低，不可回到初始状态，故 B 错误；气体等容增压，吸热，温度升高，接着等压增容，温度升高，不可回到初始状态，故 C 错误；气体等容减压，放热，温度降低，接着等压增容，吸热，温度升高，可回到初始状态，故 D 正确。

评分说明：全题 15 分，每小题 3 分。
}

\item
%% number: 16
%% paperName: 1983年全国高考第 16 题 3 分
%% typeId: 4
%% type: 实验题
%% chapter: 直线运动
%% point: 游标卡尺和螺旋测微仪
%% method: 无
%% score: 3
%% degree: 950
%% duplicateId: 0
%% body:
用螺旋测微器（千分尺）测小球直径时，可动刻度 H 的位置如右图所示。这时读出的数值是\tkanswer{13.858}，单位是\tkanswer{毫米}。
\onepicture[align=right, w=6cm]{figs/16.png}

%% answer: 13.858，毫米
\jdanswer{
13.858，毫米
}

%% memo:
\memoanswer{
固定刻度读 $13.5\Umm$，可动刻度读 $35.8\times0.01\Umm = 0.358\Umm$，故螺旋测微器读数是 $13.858\Umm$。

评分说明：本小题 3 分。13.858，毫米（1分）。数值写成 13.857 或 13.859 并且单位正确的，也给 1 分。数值或单位有错的，给 0 分。
}

\item
%% number: 17
%% paperName: 1983年全国高考第 17 题 2 分
%% typeId: 4
%% type: 实验题
%% chapter: 电路
%% point: 多用表的使用
%% method: 无
%% score: 2
%% degree: 950
%% duplicateId: 0
%% body:
用伏特表的 $0\sim3\UV$ 档测电路中某两点间的电压时，表的指针位置如图 1 所示。从伏特表读出的电压是\tkanswer{2.0}$\UV$。用安培表的 $0\sim0.6\UA$ 档测某一电路中的电流时，表的指针位置如图 2 所示。从安培表读出的电流是\tkanswer{0.48}$\UA$。

\twopicture[numA=1, numB=2, hA=2.6cm, hB=2.6cm, align=right]
{figs/17a.png}
{figs/17b.png}

%% answer: 2.0 V；0.48 A
\jdanswer{
$2.0\UV$，$0.48\UA$
}

%% memo:
\memoanswer{
两电表的分度值为 0.02，估读时使用二分法，即估读半小格。

评分说明：$2.0\UV$（1分）。答作 $2.2\UV$ 的，也给 1 分。数值或单位有错的，给 0 分。$0.48\UA$（1分）。答作 $0.480\UA$ 的，或 $0.481\UA$ 的，或 $0.482\UA$ 的，也给 1 分。数值或单位有错的，给 0 分。
}

\item
%% number: 18
%% paperName: 1983年全国高考第 18 题 3 分
%% typeId: 4
%% type: 实验题
%% chapter: 力物体的平衡
%% point: 力
%% method: 无
%% score: 3
%% degree: 750
%% duplicateId: 0
%% body:
给定一个空玻璃瓶（如图所示），要利用天平和水来测定这个玻璃瓶刻度线下的容积，测定中主要应进行。

（A）用天平称量瓶装着水时瓶和水的总质量

（B）用天平称量空瓶的质量

（C）算出瓶内水的质量，求出瓶的容积

（D）调节天平横梁两端的螺旋，使天平平衡

（E）调节天平底板下面的螺旋，使天平的底板成为水平

把以上各项的英文字母代号按实验的合理顺序填写在下面横线上空白处。

（1）\tkanswer{E}，（2）\tkanswer{D}，（3）\tkanswer{B}，（4）\tkanswer{A}，（5）\tkanswer{C}。
\onepicture[align=right, w=3.5cm]{figs/18.png}

%% answer: EDBAC
\jdanswer{
EDBAC
}

%% memo:
\memoanswer{
参考答案：EDBAC。答作 E，D，A，B，C 的，给 2 分，其他任何顺序都不得给分。

评分说明：本小题 3 分。
}

\item
%% number: 19
%% paperName: 1983年全国高考第 19 题 8 分
%% typeId: 4
%% type: 实验题
%% chapter: 电路
%% point: 输出功率极值
%% method: 无
%% score: 8
%% degree: 550
%% duplicateId: 0
%% body:
电源的输出功率 $P$ 跟外电路的电阻 $R$ 有关。图 1 是研究它们关系的实验电路。为了便于进行实验和保护蓄电池，给蓄电池串联了一个定值电阻 $R_{0}$，把它们一起看作电源（图中虚线框内部分）。于是电源的内电阻就是蓄电池的内电阻和定值电阻 $R_{0}$ 之和，用 $r$ 表示。电源的电动势用 $E$ 表示。

\begin{enumerate}
	\item
	在下面空白处列出 $P$ 跟 $E$、$r$、$R$ 的关系式（$R$ 中包含安培表的内阻；伏特表中电流不计。）

	\item
	在图 2 中，按照图 1 画出连线，把所示的器件连接成实验电路。

	下面给出六组实验数据：

	\begin{center}
	\begin{tabular}{|c|c|c|c|c|c|c|}
	\hline
	$I$（A） & 0.20 & 0.28 & 0.36 & 0.44 & 0.52 & 0.60 \\
	\hline
	$U$（V） & 3.00 & 2.60 & 2.20 & 1.80 & 1.40 & 1.00 \\
	\hline
	$U/I$（$\Omega$） & 15 & 9.3 & 6.1 & 4.1 & 2.7 & 1.7 \\
	\hline
	$UI$（W） & 0.60 & 0.73 & 0.79 & 0.79 & 0.73 & 0.60 \\
	\hline
	\end{tabular}
	\end{center}

	\item
	上表给出了六组实验数据，根据这些数据，在图 3 的直角坐标系中画出 $P-R$ 关系图线。根据图线得出的电源输出功率的最大值是\tkanswer{0.80}$\UW$，对应的外电阻的阻值是\tkanswer{5.0}$\UO$。
\end{enumerate}
\threepicture[numA=1, numB=2, numC=3, align=right, wA=5.0cm, wB=4.0cm, wC=5.0cm]
{figs/19a.png}
{figs/19b.png}
{figs/19c.png}

%% answer: （1）P=E^2R/(R+r)^2；（2）连线如图；（3）0.80 W，5.0 Ω
\jdanswer{
\begin{enumerate}
	\item
	$P = \frac{E^{2}R}{(R+r)^{2}}$

	\item
	连线如答案图所示。

	\item
	$P-R$ 关系图线如答案图所示；$0.80\UW$，$5.0\UO$。
\end{enumerate}
}

%% memo:
\memoanswer{
（1）$P=\frac{E^{2}R}{(R+r)^{2}}$（2分），本小问不要求写出推导过程。

（2）连线如下图（3分）。
\onepicture[align=right, w=8cm]{figs/19_an1.png}

只有 $R_{0}$ 未接入而其它正确的，给 2 分。安培表或伏特表或两者的正负端接错，而其他正确的，给 1 分。变阻器的接法跟下图不同，但接法正确的，不影响评分。

（3）$P-R$ 关系图线如下图（2分）。
\onepicture[align=right, w=9cm]{figs/19_an2.png}

$0.80\UW$，$5.0\UO$（1分）。作图点在正确位置附近，曲线基本平滑即可。作图点点错两个及以上的，图线给 0 分。只是作图点点错一个的，图线给 1 分。最大输出功率值允许在 $0.79\UW$ 以上到 $0.81\UW$ 之间，对应的外电阻的阻值允许在 $4.6\UO$ 到 $5.4\UO$ 之间。超出这个范围的，不给这 1 分。图线连成明显折线的，不给本小问这 3 分。
}

\item
%% number: 20
%% paperName: 1983年全国高考第 20 题 10 分
%% typeId: 6
%% type: 计算题
%% chapter: 磁场
%% point: 带电粒子在磁场中的运动
%% method: 无
%% score: 10
%% degree: 550
%% duplicateId: 0
%% body:
一个负离子，质量为 $m$，电量大小为 $q$，以速率 $v$ 垂直于屏 S 经过小孔 O 射入存在着匀强磁场的真空室中（如图）。磁感应强度 $B$ 的方向与离子的运动方向垂直，并垂直于右图纸面向里。

\begin{enumerate}
	\item
	求离子进入磁场后到达屏 S 上时的位置与 O 点的距离；

	\item
	如果离子进入磁场后经过时间 $t$ 到达位置 P，证明：直线 OP 与离子入射方向之间的夹角 $\theta$ 跟 $t$ 的关系是 $\theta = \frac{qB}{2m}t$。
\end{enumerate}
\onepicture[align=right, w=5cm]{figs/20.png}

%% answer: （1）2mv/(qB)；（2）见解析
\jdanswer{
\begin{enumerate}
	\item
	$\frac{2mv}{qB}$

	\item
	见解析。
\end{enumerate}
}

%% memo:
\memoanswer{
（1）离子的初速度与匀强磁场的方向垂直，在洛伦兹力作用下，做匀速圆周运动。设圆半径为 $r$，则
\[ qvB = m\frac{v^{2}}{r}, \]
解得 $r = \frac{mv}{qB}$。离子回到屏 S 上的位置 A 和 O 点的距离为 $AO = 2r$，解得 $AO = \frac{2mv}{qB}$。

（2）如图所示，当离子到位置 P 时，圆心角为
\[ \alpha=\frac{vt}{r}=\frac{qBt}{m}, \]
\onepicture[align=right, w=5cm]{figs/20_an.png}
由几何关系得 $\theta + \frac{180^{\circ} - \alpha}{2} = 90^{\circ}$，解得 $\alpha = 2\theta$，所以 $\theta = \frac{qB}{2m}t$。

评分说明：全题 10 分。（1）4 分，（2）6 分。（1）中，列出（a）式的，给 3 分。直接列出（b）式的，给 2 分。再利用（c）式得出正确结果的，再给 1 分。只列出（c）式的，给 0 分，单纯运算错误，扣 1 分。（2）中，得出（e）式的，给 4 分。又利用（f）式得出正确结果的，再给 2 分。只写出（d）式和（f）式，但未用（b）式的只给 3 分。有单纯运算错误的，扣 2 分（因为是证明题）。
}

\item
%% number: 21
%% paperName: 1983年全国高考第 21 题 14 分
%% typeId: 6
%% type: 计算题
%% chapter: 气体性质
%% point: 玻意耳定律
%% method: 无
%% score: 14
%% degree: 350
%% duplicateId: 0
%% body:
把上端 A 封闭、下端 B 开口的玻璃管插入水中，放掉部分空气后放手，玻璃管可以竖直地浮在水中（如下图）。设玻璃管的质量 $m = 40\Ug$，横截面积 $S = 2\Ucmq$，水面以上部分的长度 $b = 1\Ucm$，大气压强 $p_{0} = 10^{5}\UPa$。玻璃管壁厚度不计，管内空气质量不计。

\begin{enumerate}
	\item
	求玻璃管内外水面的高度差 $h$；

	\item
	用手拿住玻璃管并缓慢地把它压入水中，当管的 A 端在水面下超过某一深度时，放手后玻璃管不浮起。求这个深度；

	\item
	上一小问中，放手后玻璃管的位置是否变化？如何变化？（计算时可认为管内空气的温度不变）
\end{enumerate}
\onepicture[align=right, w=5cm]{figs/21.png}

%% answer: （1）20 cm；（2）0.52 m；（3）自行下沉
\jdanswer{
\begin{enumerate}
	\item
	$h = 20\Ucm$

	\item
	$H_{0} = 0.52\Um$

	\item
	放手后管的位置要变化，将自行下沉。
\end{enumerate}
}

%% memo:
\memoanswer{
（1）玻璃管 A 端浮在水面上方时，管受力平衡。设管中空气压强为 $p_{1}$，则管所受内外空气压力之差（竖直方向）是 $f = (p_{1} - p_{0})S$（a）。用 $\rho$ 表示水的密度，$p_{1}=p_{0}+\rho gh$（b），则 $f=\rho ghS$（c）。$f$ 应与管所受重力平衡：$\rho ghS = mg$（d），联立解得 $h=\frac{m}{\rho S}$，代入数值得 $h=\frac{40\times10^{-3}}{10^{3}\times2\times10^{-4}}\Um=20\Ucm$。

（2）管竖直没入水中后，设管 A 端的深度为 $H$，管内气柱长度为 $l$，则 A 端所在处水内压强为 $p_{A}=p_{0}+\rho gH$（f），设管内气压为 $p_{2}$，由管内水面在水下的深度可知：$p_{2}=p_{0}+\rho gH+\rho gl$（g），管所受两者压力之差（竖直方向）为 $f = (p_{2} - p_{A})S = \rho gSl$（h）。随着管的下降，管内水面也必下降，即管内水面在水下的深度增大（管内气体被压缩，压强增大，故管内液面下降）。因此，$p_{2}$ 增大，$l$ 减小，故 $f'$ 减小。当管 A 端到达某一深度 $H_{0}$ 时，$f'$ 与管所受重力相等；超过这一深度后，$f'$ 小于重力，放手后管不浮起。由此，当 $H=H_{0}$ 时，由平衡条件得 $f' = \rho gSl = mg$（i），解得 $l=\frac{m}{\rho S}$（j）。由玻意耳定律得 $p_{1}(b+h)S=p_{2}lS$，即 $\left(p_{0}+\rho gH_{0}+\rho g\frac{m}{\rho S}\right)\frac{m}{\rho S}=\left(p_{0}+\frac{m}{\rho S}\rho g\right)\left(b+\frac{m}{\rho S}\right)$，解得 $H_{0}=b\left(1+\frac{p_{0}S}{mg}\right)$，代入数值后，$H_{0}=1\times10^{-2}\left(1+\frac{10^{5}\times2\times10^{-4}}{40\times10^{-3}\times9.8}\right)\Um=0.52\Um$。

（3）由上一小问解答的分析可知，当管 A 端的深度超过 $H_{0}$ 时，$f' < mg$。故放手后管的位置要变化，将自行下沉。

评分说明：全题 14 分。（1）3 分；（2）和（3）共 11 分。（1）中，利用（a）、（b）式求出（c）式的，给 2 分。直接用阿基米德原理得出管（及管内空气）所受浮力（c）式的，同样给 2 分。利用条件（d）得出结果（e）的，再给 1 分。因单纯运算或数值计算（包括单位换算）错误而结果错误的，扣 1 分。（2）、（3），这两小问的解答中考生需要通过分析得知 $f'$ 随着管的下降而减小，从而确定放手后管不浮起的条件和管位置的变化，故两小问一起定评分说明。利用（f）、（g）得出（h）式的，给 2 分。直接求浮力而得出（h）式的，同样给 2 分。利用平衡条件得出（j）式的，再给 1 分。利用玻-马定律决定 $H_{0}$ 部分，占 3 分。分析 $f'$ 随管的下降而减小，占 4 分，不要求严格论证，能说出管下降时 $l$ 减小即可；用其他话说的，正确的也可；不作分析的不给这 4 分。说出自行下沉的，再给 1 分。因单纯运算或数值计算（包括单位换算）错误而结果错误的，扣 1 分。$g$ 值取作 $10\Umsq$ 而得出 $H_{0}=0.51\Um$ 的，同样给分。
}

\item
%% number: 22
%% paperName: 1983年全国高考第 22 题 12 分
%% typeId: 6
%% type: 计算题
%% chapter: 曲线运动
%% point: 圆周运动的应用
%% method: 临界
%% score: 12
%% degree: 350
%% duplicateId: 0
%% body:
一个光滑的圆锥体固定在水平的桌面上，其轴线沿竖直方向，母线与轴线之间的夹角 $\theta = 30^{\circ}$（如右图）。一条长度为 $l$ 的绳（质量不计），一端的位置固定在圆锥体的顶点 O 处，另一端拴着一个质量为 $m$ 的小物体（物体可看作质点，绳长小于圆锥体的母线）。物体以速率 $v$ 绕圆锥体的轴线做水平匀速圆周运动（物体和绳在上图中都没画出）。

（要求说明每问解法的根据）

\begin{enumerate}
	\item
	当 $v = \sqrt{\frac{1}{6}gl}$，求绳对物体的拉力；

	\item
	当 $v = \sqrt{\frac{3}{2}gl}$，求绳对物体的拉力。
\end{enumerate}
\onepicture[align=right, w=5cm]{figs/22.png}

%% answer: （1）1.03mg；（2）2mg
\jdanswer{
\begin{enumerate}
	\item
	$T = 1.03mg$

	\item
	$T = 2mg$
\end{enumerate}
}

%% memo:
\memoanswer{
物体做水平匀速圆周运动有两种可能：一种是物体与锥体表面接触，如图 1 所示；一种是物体与锥体表面不接触，如图 2 所示。
\onepicture[align=right, w=4.5cm]{figs/22_an1.png}

当接触时，物体受力如图 1 所示，$T$ 是绳对物体的拉力，$N$ 是支持力，$mg$ 是重力，物体与锥面间无摩擦。将力沿水平方向和竖直方向分解，由牛顿运动定律和受力分析得
\[ T\sin\theta-N\cos\theta=ma=m\frac{v^{2}}{l\sin\theta} \tag{①} \]
\[ T\cos\theta+N\sin\theta=mg \tag{②} \]
由①②两式消去 $T$，可得 $N$ 跟 $v$ 的关系为
\[ N=mg\sin\theta-m\frac{v^{2}\cos\theta}{l\sin\theta}, \]
在 $\theta$ 给定后，$v$ 越大，$N$ 就越小。当 $v=\sqrt{\frac{gl\sin^{2}\theta}{\cos\theta}}$ 时，$N=0$，令 $v_{b}$ 表示这个速率，并将 $\theta=30^{\circ}$ 代入，可得
\[ v_{b}=\sqrt{\frac{\sqrt{3}gl}{6}} \tag{③} \]
因为 $N$ 是支持力，最小等于 0，所以当 $v > v_{b}$ 时，物体不再与锥面接触。

（1）当 $v = \sqrt{\frac{1}{6}gl}$ 时，因 $v < v_{b}$，故物体与锥面接触，由①②两式消去 $N$，可得
\[ T=m\frac{v^{2}}{l}+mg\cos\theta=\frac{1+3\sqrt{3}}{6}mg, \]
或：$T=1.03mg$。

（2）当 $v = \sqrt{\frac{3}{2}gl}$ 时，因 $v > v_{b}$，故物体与锥面不接触，这时物体只受重力和绳子拉力作用（如图 2 所示）。用 $\alpha$ 表示绳与圆锥体轴线之间的夹角，将力沿水平方向和竖直方向分解，由牛顿运动定律和受力分析得
\onepicture[align=right, w=4.5cm]{figs/22_an2.png}
\[ T\sin\alpha = m\frac{v^{2}}{l\sin\alpha} \tag{④} \]
\[ T\cos\alpha = mg \tag{⑤} \]
将 $v = \sqrt{\frac{3}{2}gl}$ 代入④式，由④⑤两式消去 $\alpha$，可得
\[ 2T^{2}-3mgT-2m^{2}g^{2}=0, \]
解此方程，取合理值，得 $T=2mg$。

评分说明：全题 12 分。本题要求考生说明每问解法的根据，即要求得出（c）式，并将（1）、（2）两问中的速率与（c）式相比较，这部分内容占 6 分。不论考生用什么方法解题，得出（c）式的给 4 分，再将（1）、（2）两问中的速率与（c）式比较的，再各给 1 分。在（1）中，列（a）、（b）式及求解占 3 分；（a）、（b）两式中有一个列错的，扣 2 分；单纯运算错误，扣 1 分；答案最后结果写作 $T = mg$ 的，不扣分。在（2）中，列（d）、（e）式及求解占 3 分；（d）、（e）两式中有一个列错的，扣 2 分；单纯运算错误，扣 1 分；若误认为 $\alpha=30^{\circ}$，扣 2 分。
}

\end{enumerate}

\end{document}
